% 本文件由账本已入列事件生成；锚点首次呈现仅连接可核的年序与材料限度。
\section{辽、西夏与宋之间的同步年序}
以下补入按同年并列的政权与地区组织，而不是以一个朝廷的年号吞没其对手。每个可点击事件名保留账本 ID；叙述只将材料明确记录的动作置入年序，并把实施、规模、损失与动机留在可检验的限度内。

\subsection{960年：宋、辽与北汉／高丽（年序桥接）}
《宋史》、《辽史》、《续资治通鉴长编》为这一年留下了可回查的年序入口。

\noindent\textbf{宋、辽与北汉／高丽（年序桥接）。}\eventanchor{SY-960-CHRONOLOGY-BRIDGE}{【C级年序桥接】保留960 年的多政权并行检索入口；本条不主张所列政权在当年发生直接互动，具体事件须另以单项材料入账}。

\subsection{961年：宋、辽与北汉／高丽（年序桥接）、宋与辽}
《宋史》、《辽史》、《续资治通鉴长编》为这一年留下了可回查的年序入口。

\noindent\textbf{宋、辽与北汉／高丽（年序桥接）。}\eventanchor{SY-961-CHRONOLOGY-BRIDGE}{【C级年序桥接】保留961 年的多政权并行检索入口；本条不主张所列政权在当年发生直接互动，具体事件须另以单项材料入账}。

\noindent\textbf{宋与辽。}\eventanchor{SYD-961-LIAO-001}{乙未，李繼勳敗北漢軍，俘遼州刺史傅廷彥、弟勳來獻}。

\subsection{962年：宋、辽与北汉／高丽（年序桥接）}
《宋史》、《辽史》、《续资治通鉴长编》为这一年留下了可回查的年序入口。

\noindent\textbf{宋、辽与北汉／高丽（年序桥接）。}\eventanchor{SY-962-CHRONOLOGY-BRIDGE}{【C级年序桥接】保留962 年的多政权并行检索入口；本条不主张所列政权在当年发生直接互动，具体事件须另以单项材料入账}。

\subsection{963年：宋、辽与北汉／高丽（年序桥接）、宋与辽}
《宋史》、《辽史》、《续资治通鉴长编》为这一年留下了可回查的年序入口。

\noindent\textbf{宋、辽与北汉／高丽（年序桥接）。}\eventanchor{SY-963-CHRONOLOGY-BRIDGE}{【C级年序桥接】保留963 年的多政权并行检索入口；本条不主张所列政权在当年发生直接互动，具体事件须另以单项材料入账}。

\noindent\textbf{宋与辽。}\eventanchor{SYD-963-LIAO-002}{己亥，契丹幽州岐溝關使柴廷翰等來降}；\eventanchor{SYD-963-LIAO-086}{戊寅，北漢引契丹兵攻平晉，遣洺州防禦使郭進等救之}。

\subsection{964年：宋、辽与北汉／高丽（年序桥接）、宋与辽}
《宋史》、《辽史》、《续资治通鉴长编》为这一年留下了可回查的年序入口。

\noindent\textbf{宋、辽与北汉／高丽（年序桥接）。}\eventanchor{SY-964-CHRONOLOGY-BRIDGE}{【C级年序桥接】保留964 年的多政权并行检索入口；本条不主张所列政权在当年发生直接互动，具体事件须另以单项材料入账}。

\noindent\textbf{宋与辽。}\eventanchor{SYD-964-LIAO-003}{二月戊申朔，北漢遼州刺史杜延韜以城來降}。

\subsection{965年：宋、辽与北汉／高丽（年序桥接）}
《宋史》、《辽史》、《续资治通鉴长编》为这一年留下了可回查的年序入口。

\noindent\textbf{宋、辽与北汉／高丽（年序桥接）。}\eventanchor{SY-965-CHRONOLOGY-BRIDGE}{【C级年序桥接】保留965 年的多政权并行检索入口；本条不主张所列政权在当年发生直接互动，具体事件须另以单项材料入账}。

\subsection{966年：宋、辽与北汉／高丽（年序桥接）、宋与辽}
《宋史》、《辽史》、《续资治通鉴长编》为这一年留下了可回查的年序入口。

\noindent\textbf{宋、辽与北汉／高丽（年序桥接）。}\eventanchor{SY-966-CHRONOLOGY-BRIDGE}{【C级年序桥接】保留966 年的多政权并行检索入口；本条不主张所列政权在当年发生直接互动，具体事件须另以单项材料入账}。

\noindent\textbf{宋与辽。}\eventanchor{SYD-966-LIAO-004}{丁巳，契丹天德軍節度使于延超與其子來降}。

\subsection{967年：宋、辽与北汉／高丽（年序桥接）}
《宋史》、《辽史》、《续资治通鉴长编》为这一年留下了可回查的年序入口。

\noindent\textbf{宋、辽与北汉／高丽（年序桥接）。}\eventanchor{SY-967-CHRONOLOGY-BRIDGE}{【C级年序桥接】保留967 年的多政权并行检索入口；本条不主张所列政权在当年发生直接互动，具体事件须另以单项材料入账}。

\subsection{968年：宋、辽与北汉／高丽（年序桥接）}
《宋史》、《辽史》、《续资治通鉴长编》为这一年留下了可回查的年序入口。

\noindent\textbf{宋、辽与北汉／高丽（年序桥接）。}\eventanchor{SY-968-CHRONOLOGY-BRIDGE}{【C级年序桥接】保留968 年的多政权并行检索入口；本条不主张所列政权在当年发生直接互动，具体事件须另以单项材料入账}。

\subsection{969年：宋、辽与北汉／高丽（年序桥接）、宋与辽}
《宋史》、《辽史》、《续资治通鉴长编》为这一年留下了可回查的年序入口。

\noindent\textbf{宋、辽与北汉／高丽（年序桥接）。}\eventanchor{SY-969-CHRONOLOGY-BRIDGE}{【C级年序桥接】保留969 年的多政权并行检索入口；本条不主张所列政权在当年发生直接互动，具体事件须另以单项材料入账}。

\noindent\textbf{宋与辽。}\eventanchor{SYD-969-LIAO-005}{己未，何繼筠敗契丹於陽曲，斬首數千級，俘武州刺史王彥符以獻，命陳示所獲首級、鎧甲於城下}；\eventanchor{SYD-969-LIAO-087}{命彰德軍節度使韓仲贇為北面都部署，彰義軍節度使郭延義副之，以防契丹}。

\subsection{970年：宋、辽与北汉／高丽（年序桥接）、宋与辽}
《宋史》、《辽史》、《续资治通鉴长编》为这一年留下了可回查的年序入口。

\noindent\textbf{宋、辽与北汉／高丽（年序桥接）。}\eventanchor{SY-970-CHRONOLOGY-BRIDGE}{【C级年序桥接】保留970 年的多政权并行检索入口；本条不主张所列政权在当年发生直接互动，具体事件须另以单项材料入账}。

\noindent\textbf{宋与辽。}\eventanchor{SYD-970-LIAO-006}{癸亥，定州駐泊都監田欽祚敗契丹於遂城}。

\subsection{971年：宋、辽与北汉／高丽（年序桥接）}
《宋史》、《辽史》、《续资治通鉴长编》为这一年留下了可回查的年序入口。

\noindent\textbf{宋、辽与北汉／高丽（年序桥接）。}\eventanchor{SY-971-CHRONOLOGY-BRIDGE}{【C级年序桥接】保留971 年的多政权并行检索入口；本条不主张所列政权在当年发生直接互动，具体事件须另以单项材料入账}。

\subsection{972年：宋、辽与北汉／高丽（年序桥接）}
《宋史》、《辽史》、《续资治通鉴长编》为这一年留下了可回查的年序入口。

\noindent\textbf{宋、辽与北汉／高丽（年序桥接）。}\eventanchor{SY-972-CHRONOLOGY-BRIDGE}{【C级年序桥接】保留972 年的多政权并行检索入口；本条不主张所列政权在当年发生直接互动，具体事件须另以单项材料入账}。

\subsection{973年：宋、辽与北汉／高丽（年序桥接）}
《宋史》、《辽史》、《续资治通鉴长编》为这一年留下了可回查的年序入口。

\noindent\textbf{宋、辽与北汉／高丽（年序桥接）。}\eventanchor{SY-973-CHRONOLOGY-BRIDGE}{【C级年序桥接】保留973 年的多政权并行检索入口；本条不主张所列政权在当年发生直接互动，具体事件须另以单项材料入账}。

\subsection{974年：宋、辽与北汉／高丽（年序桥接）}
《宋史》、《辽史》、《续资治通鉴长编》为这一年留下了可回查的年序入口。

\noindent\textbf{宋、辽与北汉／高丽（年序桥接）。}\eventanchor{SY-974-CHRONOLOGY-BRIDGE}{【C级年序桥接】保留974 年的多政权并行检索入口；本条不主张所列政权在当年发生直接互动，具体事件须另以单项材料入账}。

\subsection{975年：宋、辽与北汉／高丽（年序桥接）、宋与辽}
《宋史》、《辽史》、《续资治通鉴长编》为这一年留下了可回查的年序入口。

\noindent\textbf{宋、辽与北汉／高丽（年序桥接）。}\eventanchor{SY-975-CHRONOLOGY-BRIDGE}{【C级年序桥接】保留975 年的多政权并行检索入口；本条不主张所列政权在当年发生直接互动，具体事件须另以单项材料入账}。

\noindent\textbf{宋与辽。}\eventanchor{SYD-975-LIAO-007}{庚辰，遣閣門使郝崇信、太常丞呂端使契丹}；\eventanchor{SYD-975-LIAO-088}{己亥，契丹遣使克沙骨慎思以書來講和}。

\subsection{976年：宋、辽与北汉／高丽（年序桥接）、宋与辽}
《宋史》、《辽史》、《续资治通鉴长编》为这一年留下了可回查的年序入口。

\noindent\textbf{宋、辽与北汉／高丽（年序桥接）。}\eventanchor{SY-976-CHRONOLOGY-BRIDGE}{【C级年序桥接】保留976 年的多政权并行检索入口；本条不主张所列政权在当年发生直接互动，具体事件须另以单项材料入账}。

\noindent\textbf{宋与辽。}\eventanchor{SYD-976-LIAO-008}{契丹遣使耶律延寧頁以御衣、玉帶、名馬、散馬、白鶻來賀長春節}；\eventanchor{SYD-976-LIAO-089}{己丑，遣著作郎馮正、佐郎張玘使契丹告哀}。

\subsection{977年：宋、辽与北汉／高丽（年序桥接）、宋与辽}
《宋史》、《辽史》、《续资治通鉴长编》为这一年留下了可回查的年序入口。

\noindent\textbf{宋、辽与北汉／高丽（年序桥接）。}\eventanchor{SY-977-CHRONOLOGY-BRIDGE}{【C级年序桥接】保留977 年的多政权并行检索入口；本条不主张所列政权在当年发生直接互动，具体事件须另以单项材料入账}。

\noindent\textbf{宋与辽。}\eventanchor{SYD-977-LIAO-009}{丁酉，契丹遣使來會葬}；\eventanchor{SYD-977-LIAO-090}{二月甲午，契丹遣使來賀即位及正旦}。

\subsection{978年：宋、辽与北汉／高丽（年序桥接）、宋与辽}
《宋史》、《辽史》、《续资治通鉴长编》为这一年留下了可回查的年序入口。

\noindent\textbf{宋、辽与北汉／高丽（年序桥接）。}\eventanchor{SY-978-CHRONOLOGY-BRIDGE}{【C级年序桥接】保留978 年的多政权并行检索入口；本条不主张所列政权在当年发生直接互动，具体事件须另以单项材料入账}。

\noindent\textbf{宋与辽。}\eventanchor{SYD-978-LIAO-010}{冬十月癸丑朔，契丹遣使來賀乾明節}；\eventanchor{SYD-978-LIAO-091}{癸巳，遣李從吉等使契丹}。

\subsection{979年：宋、辽与北汉／高丽（年序桥接）、宋与辽}
《宋史》、《辽史》、《续资治通鉴长编》为这一年留下了可回查的年序入口。

\noindent\textbf{宋、辽与北汉／高丽（年序桥接）。}\eventanchor{SY-979-CHRONOLOGY-BRIDGE}{【C级年序桥接】保留979 年的多政权并行检索入口；本条不主张所列政权在当年发生直接互动，具体事件须另以单项材料入账}。

\noindent\textbf{宋与辽。}\eventanchor{SYD-979-LIAO-011}{丙午，鎮州都鈐轄劉廷翰及契丹戰于遂城西，大敗之，斬首萬三百級，獲三將、馬萬匹}；\eventanchor{SYD-979-LIAO-092}{庚申，帝復自將伐契丹}。

\subsection{980年：宋、辽与北汉／高丽（年序桥接）、宋与辽}
《宋史》、《辽史》、《续资治通鉴长编》为这一年留下了可回查的年序入口。

\noindent\textbf{宋、辽与北汉／高丽（年序桥接）。}\eventanchor{SY-980-CHRONOLOGY-BRIDGE}{【C级年序桥接】保留980 年的多政权并行检索入口；本条不主张所列政权在当年发生直接互动，具体事件须另以单项材料入账}。

\noindent\textbf{宋与辽。}\eventanchor{SYD-980-LIAO-012}{關南與契丹戰，大破之}；\eventanchor{SYD-980-LIAO-093}{癸巳，代州言，宣徽南院使潘美敗契丹之師於雁門，殺其駙馬侍中蕭咄李，獲都指揮使李重誨}。

\subsection{981年：宋、辽与北汉／高丽（年序桥接）、宋与辽}
《宋史》、《辽史》、《续资治通鉴长编》为这一年留下了可回查的年序入口。

\noindent\textbf{宋、辽与北汉／高丽（年序桥接）。}\eventanchor{SY-981-CHRONOLOGY-BRIDGE}{【C级年序桥接】保留981 年的多政权并行检索入口；本条不主张所列政权在当年发生直接互动，具体事件须另以单项材料入账}。

\noindent\textbf{宋与辽。}\eventanchor{SYD-981-LIAO-013}{平塞軍與契丹戰，破之}；\eventanchor{SYD-981-LIAO-094}{七月丙午，詔渤海琰府王助討契丹}。

\subsection{982年：宋、辽与北汉／高丽（年序桥接）、宋与辽}
《宋史》、《辽史》、《续资治通鉴长编》为这一年留下了可回查的年序入口。

\noindent\textbf{宋、辽与北汉／高丽（年序桥接）。}\eventanchor{SY-982-CHRONOLOGY-BRIDGE}{【C级年序桥接】保留982 年的多政权并行检索入口；本条不主张所列政权在当年发生直接互动，具体事件须另以单项材料入账}。

\noindent\textbf{宋与辽。}\eventanchor{SYD-982-LIAO-014}{己未，府州破契丹於新澤砦，獲其將校以下百人}；\eventanchor{SYD-982-LIAO-095}{閏月戊子朔，豐州與契丹戰，破之，獲其天德軍節度使蕭太}。

\subsection{983年：宋、辽与北汉／高丽（年序桥接）、宋与辽}
《宋史》、《辽史》、《续资治通鉴长编》为这一年留下了可回查的年序入口。

\noindent\textbf{宋、辽与北汉／高丽（年序桥接）。}\eventanchor{SY-983-CHRONOLOGY-BRIDGE}{【C级年序桥接】保留983 年的多政权并行检索入口；本条不主张所列政权在当年发生直接互动，具体事件须另以单项材料入账}。

\noindent\textbf{宋与辽。}\eventanchor{SYD-983-LIAO-015}{豐州破契丹兵，降三千餘帳}。

\subsection{984年：宋、辽与北汉／高丽（年序桥接）、宋与西夏}
《宋史》、《辽史》、《续资治通鉴长编》、《金史》为这一年留下了可回查的年序入口。

\noindent\textbf{宋、辽与北汉／高丽（年序桥接）。}\eventanchor{SY-984-CHRONOLOGY-BRIDGE}{【C级年序桥接】保留984 年的多政权并行检索入口；本条不主张所列政权在当年发生直接互动，具体事件须另以单项材料入账}。

\noindent\textbf{宋与西夏。}\eventanchor{SYD-984-XIA-001}{夏州言，掩擊李繼遷，獲其母妻，俘千四百餘帳，繼遷走}。

\subsection{985年：宋、辽与北汉／高丽（年序桥接）、宋与西夏}
《宋史》、《辽史》、《续资治通鉴长编》、《金史》为这一年留下了可回查的年序入口。

\noindent\textbf{宋、辽与北汉／高丽（年序桥接）。}\eventanchor{SY-985-CHRONOLOGY-BRIDGE}{【C级年序桥接】保留985 年的多政权并行检索入口；本条不主张所列政权在当年发生直接互动，具体事件须另以单项材料入账}。

\noindent\textbf{宋与西夏。}\eventanchor{SYD-985-XIA-002}{辛丑，夏州行營破西蕃息利族，斬其代州刺史折羅遇并弟埋乞，又破保、洗兩族，降五十餘族}；\eventanchor{SYD-985-XIA-053}{乙未，夏州李繼遷誘殺汝州團練使曹光實}。

\subsection{986年：宋、辽与北汉／高丽（年序桥接）、宋与辽}
《宋史》、《辽史》、《续资治通鉴长编》为这一年留下了可回查的年序入口。

\noindent\textbf{宋、辽与北汉／高丽（年序桥接）。}\eventanchor{SY-986-CHRONOLOGY-BRIDGE}{【C级年序桥接】保留986 年的多政权并行检索入口；本条不主张所列政权在当年发生直接互动，具体事件须另以单项材料入账}。

\noindent\textbf{宋与辽。}\eventanchor{SYD-986-LIAO-016}{代州副部署盧漢贇敗契丹於土鐙堡，斬獲甚眾，殺監軍舍利二人}；\eventanchor{SYD-986-LIAO-096}{會契丹十萬眾復陷寰州，楊業護送遷民遇之，苦戰力盡，為所禽，守節而死}。

\subsection{987年：宋、辽与北汉／高丽（年序桥接）}
《宋史》、《辽史》、《续资治通鉴长编》为这一年留下了可回查的年序入口。

\noindent\textbf{宋、辽与北汉／高丽（年序桥接）。}\eventanchor{SY-987-CHRONOLOGY-BRIDGE}{【C级年序桥接】保留987 年的多政权并行检索入口；本条不主张所列政权在当年发生直接互动，具体事件须另以单项材料入账}。

\subsection{988年：宋、辽与北汉／高丽（年序桥接）、宋与辽}
《宋史》、《辽史》、《续资治通鉴长编》为这一年留下了可回查的年序入口。

\noindent\textbf{宋、辽与北汉／高丽（年序桥接）。}\eventanchor{SY-988-CHRONOLOGY-BRIDGE}{【C级年序桥接】保留988 年的多政权并行检索入口；本条不主张所列政权在当年发生直接互动，具体事件须另以单项材料入账}。

\noindent\textbf{宋与辽。}\eventanchor{SYD-988-LIAO-017}{己丑，郭守文破契丹于唐河}。

\subsection{989年：宋、辽与北汉／高丽（年序桥接）、宋与辽}
《宋史》、《辽史》、《续资治通鉴长编》为这一年留下了可回查的年序入口。

\noindent\textbf{宋、辽与北汉／高丽（年序桥接）。}\eventanchor{SY-989-CHRONOLOGY-BRIDGE}{【C级年序桥接】保留989 年的多政权并行检索入口；本条不主张所列政权在当年发生直接互动，具体事件须另以单项材料入账}。

\noindent\textbf{宋与辽。}\eventanchor{SYD-989-LIAO-018}{辛丑，契丹犯威虜軍，崇儀使尹繼倫擊破之，殺其相皮室，大將于越遁去}。

\subsection{990年：宋、辽与北汉／高丽（年序桥接）}
《宋史》、《辽史》、《续资治通鉴长编》为这一年留下了可回查的年序入口。

\noindent\textbf{宋、辽与北汉／高丽（年序桥接）。}\eventanchor{SY-990-CHRONOLOGY-BRIDGE}{【C级年序桥接】保留990 年的多政权并行检索入口；本条不主张所列政权在当年发生直接互动，具体事件须另以单项材料入账}。

\subsection{991年：宋、辽与北汉／高丽（年序桥接）、宋与辽、宋与西夏}
《宋史》、《辽史》、《续资治通鉴长编》、《金史》为这一年留下了可回查的年序入口。

\noindent\textbf{宋、辽与北汉／高丽（年序桥接）。}\eventanchor{SY-991-CHRONOLOGY-BRIDGE}{【C级年序桥接】保留991 年的多政权并行检索入口；本条不主张所列政权在当年发生直接互动，具体事件须另以单项材料入账}。

\noindent\textbf{宋与辽。}\eventanchor{SYD-991-LIAO-019}{是歲，女真表請伐契丹，詔不許，自是遂屬契丹}。

\noindent\textbf{宋与西夏。}\eventanchor{SYD-991-XIA-003}{丙子，遣商州團練使翟守素帥兵援趙保忠于夏州}。

\subsection{992年：宋、辽与北汉／高丽（年序桥接）}
《宋史》、《辽史》、《续资治通鉴长编》为这一年留下了可回查的年序入口。

\noindent\textbf{宋、辽与北汉／高丽（年序桥接）。}\eventanchor{SY-992-CHRONOLOGY-BRIDGE}{【C级年序桥接】保留992 年的多政权并行检索入口；本条不主张所列政权在当年发生直接互动，具体事件须另以单项材料入账}。

\subsection{993年：宋、辽与北汉／高丽（年序桥接）}
《宋史》、《辽史》、《续资治通鉴长编》为这一年留下了可回查的年序入口。

\noindent\textbf{宋、辽与北汉／高丽（年序桥接）。}\eventanchor{SY-993-CHRONOLOGY-BRIDGE}{【C级年序桥接】保留993 年的多政权并行检索入口；本条不主张所列政权在当年发生直接互动，具体事件须另以单项材料入账}。

\subsection{994年：宋、辽与北汉／高丽（年序桥接）、宋与辽、宋与西夏}
《宋史》、《辽史》、《续资治通鉴长编》、《金史》为这一年留下了可回查的年序入口。

\noindent\textbf{宋、辽与北汉／高丽（年序桥接）。}\eventanchor{SY-994-CHRONOLOGY-BRIDGE}{【C级年序桥接】保留994 年的多政权并行检索入口；本条不主张所列政权在当年发生直接互动，具体事件须另以单项材料入账}。

\noindent\textbf{宋与辽。}\eventanchor{SYD-994-LIAO-020}{高麗遣使，以契丹來侵乞師}。

\noindent\textbf{宋与西夏。}\eventanchor{SYD-994-XIA-004}{三月乙亥，趙保忠為趙保吉所襲，奔還夏州，指揮使趙光嗣執之以獻，李繼隆帥師入夏州}。

\subsection{995年：宋、辽与北汉／高丽（年序桥接）、宋与辽}
《宋史》、《辽史》、《续资治通鉴长编》为这一年留下了可回查的年序入口。

\noindent\textbf{宋、辽与北汉／高丽（年序桥接）。}\eventanchor{SY-995-CHRONOLOGY-BRIDGE}{【C级年序桥接】保留995 年的多政权并行检索入口；本条不主张所列政权在当年发生直接互动，具体事件须另以单项材料入账}。

\noindent\textbf{宋与辽。}\eventanchor{SYD-995-LIAO-021}{癸亥，契丹大將韓德威誘党項勒浪、嵬族自振武犯邊，永安節度使折御卿邀擊，敗之于子河汊，勒浪等乘亂反擊德威，遂殺其將突厥大尉、司徒、舍利等，獲吐渾首領一人，德威僅以身免}；\eventanchor{SYD-995-LIAO-097}{乙酉，契丹犯雄州，知州何承矩擊敗之，斬其鐵林大將一人}。

\subsection{996年：宋、辽与北汉／高丽（年序桥接）、宋与西夏}
《宋史》、《辽史》、《续资治通鉴长编》、《金史》为这一年留下了可回查的年序入口。

\noindent\textbf{宋、辽与北汉／高丽（年序桥接）。}\eventanchor{SY-996-CHRONOLOGY-BRIDGE}{【C级年序桥接】保留996 年的多政权并行检索入口；本条不主张所列政权在当年发生直接互动，具体事件须另以单项材料入账}。

\noindent\textbf{宋与西夏。}\eventanchor{SYD-996-XIA-005}{己卯，夏州、延州行營言破李繼遷於烏白池，獲未幕軍主、吃囉指揮使等二十七人，繼遷遁}。

\subsection{997年：宋、辽与北汉／高丽（年序桥接）、宋与西夏}
《宋史》、《辽史》、《续资治通鉴长编》、《金史》为这一年留下了可回查的年序入口。

\noindent\textbf{宋、辽与北汉／高丽（年序桥接）。}\eventanchor{SY-997-CHRONOLOGY-BRIDGE}{【C级年序桥接】保留997 年的多政权并行检索入口；本条不主张所列政权在当年发生直接互动，具体事件须另以单项材料入账}。

\noindent\textbf{宋与西夏。}\eventanchor{SYD-997-XIA-006}{冬十月，夏人寇靈州，合河都部署楊瓊擊走之}。

\subsection{998年：宋、辽与北汉／高丽（年序桥接）}
《宋史》、《辽史》、《续资治通鉴长编》为这一年留下了可回查的年序入口。

\noindent\textbf{宋、辽与北汉／高丽（年序桥接）。}\eventanchor{SY-998-CHRONOLOGY-BRIDGE}{【C级年序桥接】保留998 年的多政权并行检索入口；本条不主张所列政权在当年发生直接互动，具体事件须另以单项材料入账}。

\subsection{999年：宋、辽与北汉／高丽（年序桥接）、宋与辽}
《宋史》、《辽史》、《续资治通鉴长编》为这一年留下了可回查的年序入口。

\noindent\textbf{宋、辽与北汉／高丽（年序桥接）。}\eventanchor{SY-999-CHRONOLOGY-BRIDGE}{【C级年序桥接】保留999 年的多政权并行检索入口；本条不主张所列政权在当年发生直接互动，具体事件须另以单项材料入账}。

\noindent\textbf{宋与辽。}\eventanchor{SYD-999-LIAO-022}{府州言官軍入契丹五合川，拔黃太尉砦，殲其眾，焚其車帳，獲馬牛萬計}；\eventanchor{SYD-999-LIAO-098}{冀州言敗契丹兵於城南，殺千餘人，奪馬百餘匹}。

\subsection{1000年：宋、辽与北汉／高丽（年序桥接）、宋与辽}
《宋史》、《辽史》、《续资治通鉴长编》为这一年留下了可回查的年序入口。

\noindent\textbf{宋、辽与北汉／高丽（年序桥接）。}\eventanchor{SY-1000-CHRONOLOGY-BRIDGE}{【C级年序桥接】保留1000 年的多政权并行检索入口；本条不主张所列政权在当年发生直接互动，具体事件须另以单项材料入账}。

\noindent\textbf{宋与辽。}\eventanchor{SYD-1000-LIAO-023}{甲子，契丹稅木監使黃顒等率屬內附，賜冠帶}；\eventanchor{SYD-1000-LIAO-099}{九月庚辰，賜契丹降人蕭肯頭名懷忠，為右領軍衛將軍、嚴州刺史}。

\subsection{1001年：宋、辽与北汉／高丽（年序桥接）、宋与辽、宋与西夏}
《宋史》、《辽史》、《续资治通鉴长编》、《金史》为这一年留下了可回查的年序入口。

\noindent\textbf{宋、辽与北汉／高丽（年序桥接）。}\eventanchor{SY-1001-CHRONOLOGY-BRIDGE}{【C级年序桥接】保留1001 年的多政权并行检索入口；本条不主张所列政权在当年发生直接互动，具体事件须另以单项材料入账}。

\noindent\textbf{宋与辽。}\eventanchor{SYD-1001-LIAO-024}{己卯，邊臣言契丹謀入寇}；\eventanchor{SYD-1001-LIAO-100}{己未，張斌破契丹于長城口}。

\noindent\textbf{宋与西夏。}\eventanchor{SYD-1001-XIA-007}{壬午，靈州言河外砦主李瓊等以城降西夏}。

\subsection{1002年：宋、辽与北汉／高丽（年序桥接）、宋与辽}
《宋史》、《辽史》、《续资治通鉴长编》为这一年留下了可回查的年序入口。

\noindent\textbf{宋、辽与北汉／高丽（年序桥接）。}\eventanchor{SY-1002-CHRONOLOGY-BRIDGE}{【C级年序桥接】保留1002 年的多政权并行检索入口；本条不主张所列政权在当年发生直接互动，具体事件须另以单项材料入账}。

\noindent\textbf{宋与辽。}\eventanchor{SYD-1002-LIAO-025}{壬戌，契丹大林砦使王昭敏等來降}。

\subsection{1003年：宋、辽与北汉／高丽（年序桥接）、宋与辽}
《宋史》、《辽史》、《续资治通鉴长编》为这一年留下了可回查的年序入口。

\noindent\textbf{宋、辽与北汉／高丽（年序桥接）。}\eventanchor{SY-1003-CHRONOLOGY-BRIDGE}{【C级年序桥接】保留1003 年的多政权并行检索入口；本条不主张所列政权在当年发生直接互动，具体事件须另以单项材料入账}。

\noindent\textbf{宋与辽。}\eventanchor{SYD-1003-LIAO-026}{契丹來侵，戰望都縣，副都部署王繼忠陷於敵，發河東廣銳兵赴援}。

\subsection{1004年：宋与辽}
《宋史》、《辽史》为这一年留下了可回查的年序入口。

\noindent\textbf{宋与辽。}\eventanchor{SYD-1004-LIAO-027}{契丹兵至澶州北，直犯前軍西陣，其大帥撻覽耀兵出陣，俄中伏弩死}；\eventanchor{SYD-1004-LIAO-101}{保、莫州、威虜、岢嵐軍及北平砦皆擊敗契丹}。

\subsection{1005年：宋与辽、宋与西夏}
《宋史》、《辽史》、《金史》为这一年留下了可回查的年序入口。

\noindent\textbf{宋与辽。}\eventanchor{SYD-1005-LIAO-028}{丙戌，遣職方郎中韓國華等使契丹}；\eventanchor{SYD-1005-LIAO-102}{二年春正月庚戌朔，以契丹講和，大赦天下，非故鬬殺、放火、強盜、偽造符印、犯贓官典、十惡至死者，悉除之}。

\noindent\textbf{宋与西夏。}\eventanchor{SYD-1005-XIA-008}{是歲，夏州、西涼府、邛部川蠻來貢}。

\subsection{1006年：宋、辽与北汉／高丽（年序桥接）}
《宋史》、《辽史》、《续资治通鉴长编》为这一年留下了可回查的年序入口。

\noindent\textbf{宋、辽与北汉／高丽（年序桥接）。}\eventanchor{SY-1006-CHRONOLOGY-BRIDGE}{【C级年序桥接】保留1006 年的多政权并行检索入口；本条不主张所列政权在当年发生直接互动，具体事件须另以单项材料入账}。

\subsection{1007年：宋、辽与北汉／高丽（年序桥接）、宋与辽、宋与西夏}
《宋史》、《辽史》、《续资治通鉴长编》、《金史》为这一年留下了可回查的年序入口。

\noindent\textbf{宋、辽与北汉／高丽（年序桥接）。}\eventanchor{SY-1007-CHRONOLOGY-BRIDGE}{【C级年序桥接】保留1007 年的多政权并行检索入口；本条不主张所列政权在当年发生直接互动，具体事件须另以单项材料入账}。

\noindent\textbf{宋与辽。}\eventanchor{SYD-1007-LIAO-029}{十二月己亥，賜近臣、契丹錦綺綾縠等物}。

\noindent\textbf{宋与西夏。}\eventanchor{SYD-1007-XIA-009}{是歲，河西六谷、夏州、沙州、大食、占城、蒲端國、西南蕃溪峒蠻來貢}。

\subsection{1008年：宋与辽、宋与西夏}
《宋史》、《辽史》、《金史》为这一年留下了可回查的年序入口。

\noindent\textbf{宋与辽。}\eventanchor{SYD-1008-LIAO-030}{契丹使上將軍蕭智可等來賀}。

\noindent\textbf{宋与西夏。}\eventanchor{SYD-1008-XIA-010}{夏州饑，請易粟，並許之}。

\subsection{1009年：宋与辽}
《宋史》、《辽史》为这一年留下了可回查的年序入口。

\noindent\textbf{宋与辽。}\eventanchor{SYD-1009-LIAO-031}{命工部侍郎馮起為契丹國信使}；\eventanchor{SYD-1009-LIAO-103}{契丹國母蕭氏卒，輟視朝}。

\subsection{1010年：宋与辽}
《宋史》、《辽史》为这一年留下了可回查的年序入口。

\noindent\textbf{宋与辽。}\eventanchor{SYD-1010-LIAO-032}{甲子，契丹國母葬，廢朝，禁邊城樂}；\eventanchor{SYD-1010-LIAO-104}{六月庚戌，邊臣言契丹饑，來市糴，詔雄州糴粟二萬石振之}。

\subsection{1011年：宋、辽与北汉／高丽（年序桥接）}
《宋史》、《辽史》、《续资治通鉴长编》为这一年留下了可回查的年序入口。

\noindent\textbf{宋、辽与北汉／高丽（年序桥接）。}\eventanchor{SY-1011-CHRONOLOGY-BRIDGE}{【C级年序桥接】保留1011 年的多政权并行检索入口；本条不主张所列政权在当年发生直接互动，具体事件须另以单项材料入账}。

\subsection{1013年：宋、辽与北汉／高丽（年序桥接）}
《宋史》、《辽史》、《续资治通鉴长编》为这一年留下了可回查的年序入口。

\noindent\textbf{宋、辽与北汉／高丽（年序桥接）。}\eventanchor{SY-1013-CHRONOLOGY-BRIDGE}{【C级年序桥接】保留1013 年的多政权并行检索入口；本条不主张所列政权在当年发生直接互动，具体事件须另以单项材料入账}。

\subsection{1014年：宋、辽与北汉／高丽（年序桥接）、宋与西夏}
《宋史》、《辽史》、《续资治通鉴长编》、《金史》为这一年留下了可回查的年序入口。

\noindent\textbf{宋、辽与北汉／高丽（年序桥接）。}\eventanchor{SY-1014-CHRONOLOGY-BRIDGE}{【C级年序桥接】保留1014 年的多政权并行检索入口；本条不主张所列政权在当年发生直接互动，具体事件须另以单项材料入账}。

\noindent\textbf{宋与西夏。}\eventanchor{SYD-1014-XIA-011}{庚申，夏州趙德明遣使詣行闕朝貢}；\eventanchor{SYD-1014-XIA-054}{是歲，夏州、西涼府、高麗、女真來貢}。

\subsection{1015年：宋、辽与北汉／高丽（年序桥接）、宋与西夏}
《宋史》、《辽史》、《续资治通鉴长编》、《金史》为这一年留下了可回查的年序入口。

\noindent\textbf{宋、辽与北汉／高丽（年序桥接）。}\eventanchor{SY-1015-CHRONOLOGY-BRIDGE}{【C级年序桥接】保留1015 年的多政权并行检索入口；本条不主张所列政权在当年发生直接互动，具体事件须另以单项材料入账}。

\noindent\textbf{宋与西夏。}\eventanchor{SYD-1015-XIA-012}{甲寅，唃廝囉聚眾數十萬，請討平夏人以自效}。

\subsection{1016年：宋、辽与北汉／高丽（年序桥接）、宋与西夏}
《宋史》、《辽史》、《续资治通鉴长编》、《金史》为这一年留下了可回查的年序入口。

\noindent\textbf{宋、辽与北汉／高丽（年序桥接）。}\eventanchor{SY-1016-CHRONOLOGY-BRIDGE}{【C级年序桥接】保留1016 年的多政权并行检索入口；本条不主张所列政权在当年发生直接互动，具体事件须另以单项材料入账}。

\noindent\textbf{宋与西夏。}\eventanchor{SYD-1016-XIA-013}{是歲，西蕃宗哥族、邛部山后蠻、夏州、甘州來貢}；\eventanchor{SYD-1016-XIA-055}{五月乙巳，邠甯環慶部署王守斌言夏州蕃騎千五百來寇慶州，內屬蕃部擊走之}。

\subsection{1017年：宋、辽与北汉／高丽（年序桥接）、宋与辽}
《宋史》、《辽史》、《续资治通鉴长编》为这一年留下了可回查的年序入口。

\noindent\textbf{宋、辽与北汉／高丽（年序桥接）。}\eventanchor{SY-1017-CHRONOLOGY-BRIDGE}{【C级年序桥接】保留1017 年的多政权并行检索入口；本条不主张所列政权在当年发生直接互动，具体事件须另以单项材料入账}。

\noindent\textbf{宋与辽。}\eventanchor{SYD-1017-LIAO-033}{壬戌，契丹使耶律准來賀承天節}。

\subsection{1018年：宋、辽与北汉／高丽（年序桥接）}
《宋史》、《辽史》、《续资治通鉴长编》为这一年留下了可回查的年序入口。

\noindent\textbf{宋、辽与北汉／高丽（年序桥接）。}\eventanchor{SY-1018-CHRONOLOGY-BRIDGE}{【C级年序桥接】保留1018 年的多政权并行检索入口；本条不主张所列政权在当年发生直接互动，具体事件须另以单项材料入账}。

\subsection{1019年：宋、辽与北汉／高丽（年序桥接）}
《宋史》、《辽史》、《续资治通鉴长编》为这一年留下了可回查的年序入口。

\noindent\textbf{宋、辽与北汉／高丽（年序桥接）。}\eventanchor{SY-1019-CHRONOLOGY-BRIDGE}{【C级年序桥接】保留1019 年的多政权并行检索入口；本条不主张所列政权在当年发生直接互动，具体事件须另以单项材料入账}。

\subsection{1020年：宋、辽与北汉／高丽（年序桥接）}
《宋史》、《辽史》、《续资治通鉴长编》为这一年留下了可回查的年序入口。

\noindent\textbf{宋、辽与北汉／高丽（年序桥接）。}\eventanchor{SY-1020-CHRONOLOGY-BRIDGE}{【C级年序桥接】保留1020 年的多政权并行检索入口；本条不主张所列政权在当年发生直接互动，具体事件须另以单项材料入账}。

\subsection{1021年：宋、辽与北汉／高丽（年序桥接）}
《宋史》、《辽史》、《续资治通鉴长编》为这一年留下了可回查的年序入口。

\noindent\textbf{宋、辽与北汉／高丽（年序桥接）。}\eventanchor{SY-1021-CHRONOLOGY-BRIDGE}{【C级年序桥接】保留1021 年的多政权并行检索入口；本条不主张所列政权在当年发生直接互动，具体事件须另以单项材料入账}。

\subsection{1022年：宋与辽}
《宋史》、《辽史》为这一年留下了可回查的年序入口。

\noindent\textbf{宋与辽。}\eventanchor{SYD-1022-LIAO-034}{八月壬寅，遣使賀契丹主及其妻生日、正旦}；\eventanchor{SYD-1022-LIAO-105}{丙寅，遣使以先帝遺留物遺契丹}。

\subsection{1023年：宋与辽}
《宋史》、《辽史》为这一年留下了可回查的年序入口。

\noindent\textbf{宋与辽。}\eventanchor{SYD-1023-LIAO-035}{庚午，契丹使初來賀長寧節}。

\subsection{1024年：宋、辽与北汉／高丽（年序桥接）}
《宋史》、《辽史》、《续资治通鉴长编》为这一年留下了可回查的年序入口。

\noindent\textbf{宋、辽与北汉／高丽（年序桥接）。}\eventanchor{SY-1024-CHRONOLOGY-BRIDGE}{【C级年序桥接】保留1024 年的多政权并行检索入口；本条不主张所列政权在当年发生直接互动，具体事件须另以单项材料入账}。

\subsection{1025年：宋、辽与北汉／高丽（年序桥接）}
《宋史》、《辽史》、《续资治通鉴长编》为这一年留下了可回查的年序入口。

\noindent\textbf{宋、辽与北汉／高丽（年序桥接）。}\eventanchor{SY-1025-CHRONOLOGY-BRIDGE}{【C级年序桥接】保留1025 年的多政权并行检索入口；本条不主张所列政权在当年发生直接互动，具体事件须另以单项材料入账}。

\subsection{1026年：宋、辽与北汉／高丽（年序桥接）}
《宋史》、《辽史》、《续资治通鉴长编》为这一年留下了可回查的年序入口。

\noindent\textbf{宋、辽与北汉／高丽（年序桥接）。}\eventanchor{SY-1026-CHRONOLOGY-BRIDGE}{【C级年序桥接】保留1026 年的多政权并行检索入口；本条不主张所列政权在当年发生直接互动，具体事件须另以单项材料入账}。

\subsection{1027年：宋、辽与北汉／高丽（年序桥接）}
《宋史》、《辽史》、《续资治通鉴长编》为这一年留下了可回查的年序入口。

\noindent\textbf{宋、辽与北汉／高丽（年序桥接）。}\eventanchor{SY-1027-CHRONOLOGY-BRIDGE}{【C级年序桥接】保留1027 年的多政权并行检索入口；本条不主张所列政权在当年发生直接互动，具体事件须另以单项材料入账}。

\subsection{1028年：宋、辽与北汉／高丽（年序桥接）}
《宋史》、《辽史》、《续资治通鉴长编》为这一年留下了可回查的年序入口。

\noindent\textbf{宋、辽与北汉／高丽（年序桥接）。}\eventanchor{SY-1028-CHRONOLOGY-BRIDGE}{【C级年序桥接】保留1028 年的多政权并行检索入口；本条不主张所列政权在当年发生直接互动，具体事件须另以单项材料入账}。

\subsection{1029年：宋、辽与北汉／高丽（年序桥接）、宋与辽}
《宋史》、《辽史》、《续资治通鉴长编》为这一年留下了可回查的年序入口。

\noindent\textbf{宋、辽与北汉／高丽（年序桥接）。}\eventanchor{SY-1029-CHRONOLOGY-BRIDGE}{【C级年序桥接】保留1029 年的多政权并行检索入口；本条不主张所列政权在当年发生直接互动，具体事件须另以单项材料入账}。

\noindent\textbf{宋与辽。}\eventanchor{SYD-1029-LIAO-036}{辛巳，詔契丹饑民所過給米，分送唐、鄧等州，以閒田處之}。

\subsection{1030年：宋、辽与北汉／高丽（年序桥接）}
《宋史》、《辽史》、《续资治通鉴长编》为这一年留下了可回查的年序入口。

\noindent\textbf{宋、辽与北汉／高丽（年序桥接）。}\eventanchor{SY-1030-CHRONOLOGY-BRIDGE}{【C级年序桥接】保留1030 年的多政权并行检索入口；本条不主张所列政权在当年发生直接互动，具体事件须另以单项材料入账}。

\subsection{1031年：宋与辽}
《宋史》、《辽史》为这一年留下了可回查的年序入口。

\noindent\textbf{宋与辽。}\eventanchor{SYD-1031-LIAO-037}{秋七月丙午朔，契丹使來告其主隆緒殂，遣使祭奠吊慰，及賀宗真立}；\eventanchor{SYD-1031-LIAO-106}{是歲，契丹主及其國母遣使來致遺留物及謝弔祭}。

\subsection{1033年：宋与辽}
《宋史》、《辽史》为这一年留下了可回查的年序入口。

\noindent\textbf{宋与辽。}\eventanchor{SYD-1033-LIAO-038}{八月甲午朔，契丹使來吊慰祭奠}。

\subsection{1035年：宋、辽与北汉／高丽（年序桥接）}
《宋史》、《辽史》、《续资治通鉴长编》为这一年留下了可回查的年序入口。

\noindent\textbf{宋、辽与北汉／高丽（年序桥接）。}\eventanchor{SY-1035-CHRONOLOGY-BRIDGE}{【C级年序桥接】保留1035 年的多政权并行检索入口；本条不主张所列政权在当年发生直接互动，具体事件须另以单项材料入账}。

\subsection{1037年：宋、辽与北汉／高丽（年序桥接）}
《宋史》、《辽史》、《续资治通鉴长编》为这一年留下了可回查的年序入口。

\noindent\textbf{宋、辽与北汉／高丽（年序桥接）。}\eventanchor{SY-1037-CHRONOLOGY-BRIDGE}{【C级年序桥接】保留1037 年的多政权并行检索入口；本条不主张所列政权在当年发生直接互动，具体事件须另以单项材料入账}。

\subsection{1039年：宋、辽与西夏（年序桥接）}
《宋史》、《辽史》、《续资治通鉴长编》为这一年留下了可回查的年序入口。

\noindent\textbf{宋、辽与西夏（年序桥接）。}\eventanchor{SY-1039-CHRONOLOGY-BRIDGE}{【C级年序桥接】保留1039 年的多政权并行检索入口；本条不主张所列政权在当年发生直接互动，具体事件须另以单项材料入账}。

\subsection{1040年：宋与辽、宋与西夏}
《宋史》、《辽史》、《金史》为这一年留下了可回查的年序入口。

\noindent\textbf{宋与辽。}\eventanchor{SYD-1040-LIAO-039}{秋七月乙丑，遣使以討元昊告契丹}；\eventanchor{SYD-1040-LIAO-107}{乙未，契丹國母復遣使來賀乾元節}。

\noindent\textbf{宋与西夏。}\eventanchor{SYD-1040-XIA-014}{賜京師、河北、陝西、河東諸軍緡錢，蠲陝西夏稅十之二，減河東所科粟}。

\subsection{1042年：宋与辽}
《宋史》、《辽史》为这一年留下了可回查的年序入口。

\noindent\textbf{宋与辽。}\eventanchor{SYD-1042-LIAO-040}{丙寅，契丹遣使來再致誓書，報徹兵}；\eventanchor{SYD-1042-LIAO-108}{己巳，契丹遣蕭英、劉六符來致書求割地}。

\subsection{1043年：宋与西夏}
《宋史》、《辽史》、《金史》为这一年留下了可回查的年序入口。

\noindent\textbf{宋与西夏。}\eventanchor{SYD-1043-XIA-015}{癸卯，遣保安軍判官邵良佐使元昊，許封冊為夏國主，歲賜絹十萬匹、茶三萬斤}；\eventanchor{SYD-1043-XIA-056}{癸巳，元昊自名曩霄，遣人來納款，稱夏國}。

\subsection{1044年：宋与辽、宋与西夏}
《宋史》、《辽史》、《金史》为这一年留下了可回查的年序入口。

\noindent\textbf{宋与辽。}\eventanchor{SYD-1044-LIAO-041}{癸未，契丹遣使來告伐夏國}；\eventanchor{SYD-1044-LIAO-109}{戊戌，命右正言余靖報使契丹}。

\noindent\textbf{宋与西夏。}\eventanchor{SYD-1044-XIA-016}{乙未，封曩霄為夏國主}。

\subsection{1045年：宋与辽}
《宋史》、《辽史》为这一年留下了可回查的年序入口。

\noindent\textbf{宋与辽。}\eventanchor{SYD-1045-LIAO-042}{丙子，契丹遣使來告伐夏國還}；\eventanchor{SYD-1045-LIAO-110}{冬十月乙卯，契丹遣使來獻九龍車及所獲夏國羊馬}。

\subsection{1046年：宋、辽与西夏（年序桥接）、宋与西夏}
《宋史》、《辽史》、《续资治通鉴长编》、《金史》为这一年留下了可回查的年序入口。

\noindent\textbf{宋、辽与西夏（年序桥接）。}\eventanchor{SY-1046-CHRONOLOGY-BRIDGE}{【C级年序桥接】保留1046 年的多政权并行检索入口；本条不主张所列政权在当年发生直接互动，具体事件须另以单项材料入账}。

\noindent\textbf{宋与西夏。}\eventanchor{SYD-1046-XIA-017}{十一月己卯，遣官議夏國公封界}。

\subsection{1048年：宋与西夏}
《宋史》、《辽史》、《金史》为这一年留下了可回查的年序入口。

\noindent\textbf{宋与西夏。}\eventanchor{SYD-1048-XIA-018}{夏四月己巳朔，封曩霄子諒祚為夏國主}。

\subsection{1049年：宋与辽}
《宋史》、《辽史》为这一年留下了可回查的年序入口。

\noindent\textbf{宋与辽。}\eventanchor{SYD-1049-LIAO-043}{己未，契丹遣使來告伐夏國}。

\subsection{1050年：宋、辽与西夏（年序桥接）、宋与辽}
《宋史》、《辽史》、《续资治通鉴长编》为这一年留下了可回查的年序入口。

\noindent\textbf{宋、辽与西夏（年序桥接）。}\eventanchor{SY-1050-CHRONOLOGY-BRIDGE}{【C级年序桥接】保留1050 年的多政权并行检索入口；本条不主张所列政权在当年发生直接互动，具体事件须另以单项材料入账}。

\noindent\textbf{宋与辽。}\eventanchor{SYD-1050-LIAO-044}{庚子，契丹遣使以伐夏師還來告}。

\subsection{1051年：宋、辽与西夏（年序桥接）}
《宋史》、《辽史》、《续资治通鉴长编》为这一年留下了可回查的年序入口。

\noindent\textbf{宋、辽与西夏（年序桥接）。}\eventanchor{SY-1051-CHRONOLOGY-BRIDGE}{【C级年序桥接】保留1051 年的多政权并行检索入口；本条不主张所列政权在当年发生直接互动，具体事件须另以单项材料入账}。

\subsection{1054年：宋、辽与西夏（年序桥接）、宋与辽}
《宋史》、《辽史》、《续资治通鉴长编》为这一年留下了可回查的年序入口。

\noindent\textbf{宋、辽与西夏（年序桥接）。}\eventanchor{SY-1054-CHRONOLOGY-BRIDGE}{【C级年序桥接】保留1054 年的多政权并行检索入口；本条不主张所列政权在当年发生直接互动，具体事件须另以单项材料入账}。

\noindent\textbf{宋与辽。}\eventanchor{SYD-1054-LIAO-045}{九月乙亥，契丹遣使來告夏國平}；\eventanchor{SYD-1054-LIAO-111}{辛巳，遣三司使王拱辰報使契丹}。

\subsection{1055年：宋与辽}
《宋史》、《辽史》为这一年留下了可回查的年序入口。

\noindent\textbf{宋与辽。}\eventanchor{SYD-1055-LIAO-046}{庚子，契丹遣使致其主宗真遺留物及謝弔祭}；\eventanchor{SYD-1055-LIAO-112}{九月戊午，契丹使來告其國主宗真殂，帝為發哀，成服於內東門幕次，遣使祭奠、吊慰及賀其子洪基立}。

\subsection{1056年：宋、辽与西夏（年序桥接）}
《宋史》、《辽史》、《续资治通鉴长编》为这一年留下了可回查的年序入口。

\noindent\textbf{宋、辽与西夏（年序桥接）。}\eventanchor{SY-1056-CHRONOLOGY-BRIDGE}{【C级年序桥接】保留1056 年的多政权并行检索入口；本条不主张所列政权在当年发生直接互动，具体事件须另以单项材料入账}。

\subsection{1057年：宋与辽、宋与西夏}
《宋史》、《辽史》、《金史》为这一年留下了可回查的年序入口。

\noindent\textbf{宋与辽。}\eventanchor{SYD-1057-LIAO-047}{冬十月乙巳，遣胡宿報使契丹}；\eventanchor{SYD-1057-LIAO-113}{九月庚子，契丹再使蕭扈、吳湛來求御容}。

\noindent\textbf{宋与西夏。}\eventanchor{SYD-1057-XIA-019}{六月戊午，夏國主諒祚遣人來謝使弔祭}。

\subsection{1058年：宋、辽与西夏（年序桥接）、宋与辽}
《宋史》、《辽史》、《续资治通鉴长编》为这一年留下了可回查的年序入口。

\noindent\textbf{宋、辽与西夏（年序桥接）。}\eventanchor{SY-1058-CHRONOLOGY-BRIDGE}{【C级年序桥接】保留1058 年的多政权并行检索入口；本条不主张所列政权在当年发生直接互动，具体事件须另以单项材料入账}。

\noindent\textbf{宋与辽。}\eventanchor{SYD-1058-LIAO-048}{二月癸卯，契丹使來告其祖母哀，輟視朝七日，遣使祭奠吊慰}；\eventanchor{SYD-1058-LIAO-114}{己-\{丑\}-，契丹遣使來謝}。

\subsection{1059年：宋、辽与西夏（年序桥接）}
《宋史》、《辽史》、《续资治通鉴长编》为这一年留下了可回查的年序入口。

\noindent\textbf{宋、辽与西夏（年序桥接）。}\eventanchor{SY-1059-CHRONOLOGY-BRIDGE}{【C级年序桥接】保留1059 年的多政权并行检索入口；本条不主张所列政权在当年发生直接互动，具体事件须另以单项材料入账}。

\subsection{1060年：宋、辽与西夏（年序桥接）}
《宋史》、《辽史》、《续资治通鉴长编》为这一年留下了可回查的年序入口。

\noindent\textbf{宋、辽与西夏（年序桥接）。}\eventanchor{SY-1060-CHRONOLOGY-BRIDGE}{【C级年序桥接】保留1060 年的多政权并行检索入口；本条不主张所列政权在当年发生直接互动，具体事件须另以单项材料入账}。

\subsection{1061年：宋、辽与西夏（年序桥接）}
《宋史》、《辽史》、《续资治通鉴长编》为这一年留下了可回查的年序入口。

\noindent\textbf{宋、辽与西夏（年序桥接）。}\eventanchor{SY-1061-CHRONOLOGY-BRIDGE}{【C级年序桥接】保留1061 年的多政权并行检索入口；本条不主张所列政权在当年发生直接互动，具体事件须另以单项材料入账}。

\subsection{1062年：宋、辽与西夏（年序桥接）、宋与西夏}
《宋史》、《辽史》、《续资治通鉴长编》、《金史》为这一年留下了可回查的年序入口。

\noindent\textbf{宋、辽与西夏（年序桥接）。}\eventanchor{SY-1062-CHRONOLOGY-BRIDGE}{【C级年序桥接】保留1062 年的多政权并行检索入口；本条不主张所列政权在当年发生直接互动，具体事件须另以单项材料入账}。

\noindent\textbf{宋与西夏。}\eventanchor{SYD-1062-XIA-020}{己-\{丑\}-，夏國主諒祚進馬，求賜書，詔賜《九經》，還其馬}。

\subsection{1063年：宋与辽}
《宋史》、《辽史》为这一年留下了可回查的年序入口。

\noindent\textbf{宋与辽。}\eventanchor{SYD-1063-LIAO-049}{六月辛卯，契丹遣蕭福延等來祭吊}；\eventanchor{SYD-1063-LIAO-115}{遣韓贄等告即位於契丹}。

\subsection{1064年：宋与辽、宋与西夏}
《宋史》、《辽史》、《金史》为这一年留下了可回查的年序入口。

\noindent\textbf{宋与辽。}\eventanchor{SYD-1064-LIAO-050}{丙辰，契丹遣耶律烈等來賀壽聖節，蕭禧等來賀明年正旦}；\eventanchor{SYD-1064-LIAO-116}{乙卯，遣兵部員外郎呂誨等四人充賀契丹太后生辰、正旦使，刑部郎中章岷等四人充賀契丹主生辰、正旦使}。

\noindent\textbf{宋与西夏。}\eventanchor{SYD-1064-XIA-021}{庚午，詔夏國精擇使人，戒勵毋紊彝章}。

\subsection{1066年：宋、辽与西夏（年序桥接）、宋与辽、宋与西夏}
《宋史》、《辽史》、《续资治通鉴长编》、《金史》为这一年留下了可回查的年序入口。

\noindent\textbf{宋、辽与西夏（年序桥接）。}\eventanchor{SY-1066-CHRONOLOGY-BRIDGE}{【C级年序桥接】保留1066 年的多政权并行检索入口；本条不主张所列政权在当年发生直接互动，具体事件须另以单项材料入账}。

\noindent\textbf{宋与辽。}\eventanchor{SYD-1066-LIAO-051}{八月庚子，遣傅卞等賀遼主生辰，張師顏等賀正旦}；\eventanchor{SYD-1066-LIAO-117}{癸酉，契丹改國號爲遼}。

\noindent\textbf{宋与西夏。}\eventanchor{SYD-1066-XIA-022}{是歲，遣使以違約數寇責夏國，諒詐獻方物謝罪}。

\subsection{1067年：宋与辽、宋与西夏}
《宋史》、《辽史》、《金史》为这一年留下了可回查的年序入口。

\noindent\textbf{宋与辽。}\eventanchor{SYD-1067-LIAO-052}{己巳，遼遣蕭傑等來賀正旦}；\eventanchor{SYD-1067-LIAO-118}{甲午，遼遣耶律好謀等來賀即位}。

\noindent\textbf{宋与西夏。}\eventanchor{SYD-1067-XIA-023}{甲申，夏國主諒祚遣使謝罪}。

\subsection{1068年：宋与辽、宋与西夏}
《宋史》、《辽史》、《金史》为这一年留下了可回查的年序入口。

\noindent\textbf{宋与辽。}\eventanchor{SYD-1068-LIAO-053}{丁卯，遣張宗益等賀遼主生辰、正旦}；\eventanchor{SYD-1068-LIAO-119}{甲子，遼遣耶律公質等來賀正旦}。

\noindent\textbf{宋与西夏。}\eventanchor{SYD-1068-XIA-024}{庚戌，賜夏國主秉常詔，許納塞門、安遠二砦歸其綏州}；\eventanchor{SYD-1068-XIA-057}{三月庚辰，夏主諒詐卒，遣使來告哀}。

\subsection{1069年：宋与辽、宋与西夏}
《宋史》、《辽史》、《金史》为这一年留下了可回查的年序入口。

\noindent\textbf{宋与辽。}\eventanchor{SYD-1069-LIAO-054}{丁-\{丑\}-，遣孫固等賀遼主生辰、正旦}；\eventanchor{SYD-1069-LIAO-120}{壬寅，遼遣耶律昌等來賀同天節}。

\noindent\textbf{宋与西夏。}\eventanchor{SYD-1069-XIA-025}{逵遣趙禼交夏人所納安遠、塞門二砦，就定地界}；\eventanchor{SYD-1069-XIA-058}{夏國請從舊蕃儀，詔許之}；\eventanchor{SYD-1069-XIA-076}{夏人渝初盟，禼請城綏州，不以易二砦，因改名綏德城}。

\subsection{1070年：宋与辽、宋与西夏}
《宋史》、《辽史》、《金史》为这一年留下了可回查的年序入口。

\noindent\textbf{宋与辽。}\eventanchor{SYD-1070-LIAO-055}{丙寅，遼遣耶律寬來賀同天節}；\eventanchor{SYD-1070-LIAO-121}{遣張景憲等賀遼主生辰、正旦}。

\noindent\textbf{宋与西夏。}\eventanchor{SYD-1070-XIA-026}{冬十月辛酉，詔延州毋納夏使}；\eventanchor{SYD-1070-XIA-059}{庚午，夏人寇鎮戎軍三川砦，巡檢趙普伏兵邀擊，敗之}；\eventanchor{SYD-1070-XIA-077}{己卯，夏人犯大順城，知慶州李復圭以方略授環慶路鈐轄李信、慶州東路都巡檢劉甫、監押種詠出戰，兵少取敗}；\eventanchor{SYD-1070-XIA-089}{甲辰，夏人寇大順城，都監燕達等擊走之}；\eventanchor{SYD-1070-XIA-095}{是月，慶州巡檢姚兕敗夏人于荔原堡}。

\subsection{1071年：宋与辽、宋与西夏}
《宋史》、《辽史》、《金史》为这一年留下了可回查的年序入口。

\noindent\textbf{宋与辽。}\eventanchor{SYD-1071-LIAO-056}{丙子，遼遣耶律紀等來賀正旦}；\eventanchor{SYD-1071-LIAO-122}{癸酉，遣楚建中等賀遼主生辰、正旦}。

\noindent\textbf{宋与西夏。}\eventanchor{SYD-1071-XIA-027}{己-\{丑\}-，種諤襲夏兵於囉兀北，大敗之，遂城囉兀}；\eventanchor{SYD-1071-XIA-060}{壬戌，遣環慶都鈐轄幵贇以兵屯邠、涇、河中，以備西夏}；\eventanchor{SYD-1071-XIA-078}{自是夏人日聚兵爲報復計，言者以諤爲稔邊患不便}。

\subsection{1072年：宋与辽}
《宋史》、《辽史》为这一年留下了可回查的年序入口。

\noindent\textbf{宋与辽。}\eventanchor{SYD-1072-LIAO-057}{癸巳，遣崔台符等賀遼主生辰、正旦}；\eventanchor{SYD-1072-LIAO-123}{己亥，遼遣蕭瑜等來賀正旦}。

\subsection{1073年：宋、辽与西夏（年序桥接）、宋与辽、宋与西夏}
《宋史》、《辽史》、《续资治通鉴长编》、《金史》为这一年留下了可回查的年序入口。

\noindent\textbf{宋、辽与西夏（年序桥接）。}\eventanchor{SY-1073-CHRONOLOGY-BRIDGE}{【C级年序桥接】保留1073 年的多政权并行检索入口；本条不主张所列政权在当年发生直接互动，具体事件须另以单项材料入账}。

\noindent\textbf{宋与辽。}\eventanchor{SYD-1073-LIAO-058}{八月壬申朔，遣賈昌衡等賀遼主生辰、正旦}；\eventanchor{SYD-1073-LIAO-124}{丙申，遼遣耶律洞等來賀正旦}。

\noindent\textbf{宋与西夏。}\eventanchor{SYD-1073-XIA-028}{二月辛卯，夏人寇秦州，都巡檢使劉惟吉敗之}。

\subsection{1074年：宋与辽}
《宋史》、《辽史》为这一年留下了可回查的年序入口。

\noindent\textbf{宋与辽。}\eventanchor{SYD-1074-LIAO-059}{丙辰，遼遣林牙蕭禧來言河東疆界，命太常少卿劉忱議之}；\eventanchor{SYD-1074-LIAO-125}{己-\{丑\}-，遼遣耶律甯等來賀正旦}。

\subsection{1075年：宋与辽}
《宋史》、《辽史》为这一年留下了可回查的年序入口。

\noindent\textbf{宋与辽。}\eventanchor{SYD-1075-LIAO-060}{丙申，遣謝景溫等賀遼主生辰、正旦}；\eventanchor{SYD-1075-LIAO-126}{丁卯，遼遣耶律景熙等來賀同天節}。

\subsection{1076年：宋与辽}
《宋史》、《辽史》为这一年留下了可回查的年序入口。

\noindent\textbf{宋与辽。}\eventanchor{SYD-1076-LIAO-061}{丙午，遣王克臣等吊慰於遼}；\eventanchor{SYD-1076-LIAO-127}{丁未，遼遣耶律運等來賀正旦}。

\subsection{1077年：宋、辽与西夏（年序桥接）、宋与辽}
《宋史》、《辽史》、《续资治通鉴长编》为这一年留下了可回查的年序入口。

\noindent\textbf{宋、辽与西夏（年序桥接）。}\eventanchor{SY-1077-CHRONOLOGY-BRIDGE}{【C级年序桥接】保留1077 年的多政权并行检索入口；本条不主张所列政权在当年发生直接互动，具体事件须另以单项材料入账}。

\noindent\textbf{宋与辽。}\eventanchor{SYD-1077-LIAO-062}{遣蘇頌等賀遼主生辰、正旦}；\eventanchor{SYD-1077-LIAO-128}{辛-\{丑\}-，遼遣耶律孝淳等來賀正旦}。

\subsection{1078年：宋、辽与西夏（年序桥接）、宋与辽}
《宋史》、《辽史》、《续资治通鉴长编》为这一年留下了可回查的年序入口。

\noindent\textbf{宋、辽与西夏（年序桥接）。}\eventanchor{SY-1078-CHRONOLOGY-BRIDGE}{【C级年序桥接】保留1078 年的多政权并行检索入口；本条不主张所列政权在当年发生直接互动，具体事件须另以单项材料入账}。

\noindent\textbf{宋与辽。}\eventanchor{SYD-1078-LIAO-063}{丙寅，遼遣耶律隆等來賀正旦}；\eventanchor{SYD-1078-LIAO-129}{甲寅，遣黃履等賀遼主生辰、正旦}。

\subsection{1079年：宋与辽、宋与西夏}
《宋史》、《辽史》、《金史》为这一年留下了可回查的年序入口。

\noindent\textbf{宋与辽。}\eventanchor{SYD-1079-LIAO-064}{庚申，遼遣蕭甯等來賀正旦}；\eventanchor{SYD-1079-LIAO-130}{甲辰，遼遣蕭晟等來賀同天節}。

\noindent\textbf{宋与西夏。}\eventanchor{SYD-1079-XIA-029}{八月丙申朔，夏人寇綏德城，都監李浦敗之}。

\subsection{1080年：宋与辽}
《宋史》、《辽史》为这一年留下了可回查的年序入口。

\noindent\textbf{宋与辽。}\eventanchor{SYD-1080-LIAO-065}{癸-\{丑\}-，遣王存等賀遼主生辰、正旦}；\eventanchor{SYD-1080-LIAO-131}{己亥，遼遣耶律永芳等來賀同天節}。

\subsection{1081年：宋与辽、宋与西夏}
《宋史》、《辽史》、《金史》为这一年留下了可回查的年序入口。

\noindent\textbf{宋与辽。}\eventanchor{SYD-1081-LIAO-066}{戊寅，遼遣蕭福全等來賀正旦}；\eventanchor{SYD-1081-LIAO-132}{夏四月癸亥，遼遣耶律祐等來賀同天節}。

\noindent\textbf{宋与西夏。}\eventanchor{SYD-1081-XIA-030}{-\{种\}-諤遣曲珍等領兵通黑水安定堡，路遇夏人，與戰，破之，斬獲甚衆}；\eventanchor{SYD-1081-XIA-061}{丙午，高遵裕以師還，夏人來追，遂潰}；\eventanchor{SYD-1081-XIA-079}{丙午，詔諭夏主左右并嵬名部族諸部首領，并許自歸}；\eventanchor{SYD-1081-XIA-090}{丁-\{丑\}-，曲珍與夏人戰於蒲桃山，敗之}；\eventanchor{SYD-1081-XIA-096}{丁-\{丑\}-，熙河經制李憲敗夏人於西市新城，獲酋首三人、首領二十余人}；\eventanchor{SYD-1081-XIA-100}{丁亥，諸軍合攻靈州，-\{种\}-諤敗夏人于黑水}；\eventanchor{SYD-1081-XIA-104}{庚申，熙河兵至女遮谷，與夏人遇，戰敗之}；\eventanchor{SYD-1081-XIA-108}{庚戌，夏兵救米脂砦，鄜延經略副使-\{种\}-諤率衆擊破之}；\eventanchor{SYD-1081-XIA-110}{庚寅，西邊守臣言夏人囚其主秉常，詔陝西、河東路討之}。

\subsection{1082年：宋与辽、宋与西夏}
《宋史》、《辽史》、《金史》为这一年留下了可回查的年序入口。

\noindent\textbf{宋与辽。}\eventanchor{SYD-1082-LIAO-067}{丁巳，遼遣耶律永端等來賀同天節}；\eventanchor{SYD-1082-LIAO-133}{壬申，遼遣耶律儀等來賀正旦}。

\noindent\textbf{宋与西夏。}\eventanchor{SYD-1082-XIA-031}{九月丁亥，夏人三十萬衆寇永樂，曲珍戰不利，裨將寇偉等死之，夏人遂圍城}；\eventanchor{SYD-1082-XIA-062}{六月辛亥朔，環慶經略司遣將與夏人戰，破之，斬其統軍嵬名妹精嵬、副統軍訛勃遇}；\eventanchor{SYD-1082-XIA-080}{壬寅，鄜延路副總管曲珍敗夏人于金湯}；\eventanchor{SYD-1082-XIA-091}{戊寅，曲珍等敗夏人於明堂川}。

\subsection{1083年：宋、辽与西夏（年序桥接）、宋与辽、宋与西夏}
《宋史》、《辽史》、《续资治通鉴长编》、《金史》为这一年留下了可回查的年序入口。

\noindent\textbf{宋、辽与西夏（年序桥接）。}\eventanchor{SY-1083-CHRONOLOGY-BRIDGE}{【C级年序桥接】保留1083 年的多政权并行检索入口；本条不主张所列政权在当年发生直接互动，具体事件须另以单项材料入账}。

\noindent\textbf{宋与辽。}\eventanchor{SYD-1083-LIAO-068}{辛亥，遼遣蕭固等來賀同天節}；\eventanchor{SYD-1083-LIAO-134}{乙酉，遣蔡京等賀遼主生辰、正旦}。

\noindent\textbf{宋与西夏。}\eventanchor{SYD-1083-XIA-032}{丙申，河東將薛義敗夏人于葭蘆西嶺}；\eventanchor{SYD-1083-XIA-063}{二月丁未，夏人數十萬衆攻蘭州，鈐轄王文鬱率死士七百餘人擊走之}；\eventanchor{SYD-1083-XIA-081}{己亥，河東將高永翼敗夏人於真卿流部}；\eventanchor{SYD-1083-XIA-092}{甲午，夏人寇蘭州，右侍禁韋定死之}；\eventanchor{SYD-1083-XIA-097}{麟、府州將郭忠詔等敗夏人於乜離抑部，詔行賞有差}；\eventanchor{SYD-1083-XIA-101}{閏月乙亥朔，夏主秉常請修貢，許之}；\eventanchor{SYD-1083-XIA-105}{三月辛卯，夏人寇蘭州，副總管李浩以衛城有功，復隴州團練使}；\eventanchor{SYD-1083-XIA-109}{是月，夏人寇麟州，知州訾虎敗之}。

\subsection{1084年：宋与辽、宋与西夏}
《宋史》、《辽史》、《金史》为这一年留下了可回查的年序入口。

\noindent\textbf{宋与辽。}\eventanchor{SYD-1084-LIAO-069}{辛卯，遼遣耶律襄等來賀正旦}；\eventanchor{SYD-1084-LIAO-135}{辛巳，遣陳睦等賀遼主生辰、正旦}。

\noindent\textbf{宋与西夏。}\eventanchor{SYD-1084-XIA-033}{冬十月乙亥，夏人寇熙河}；\eventanchor{SYD-1084-XIA-064}{癸-\{丑\}-，夏人寇蘭州，李憲等擊走之}；\eventanchor{SYD-1084-XIA-082}{癸巳，夏人寇延州安塞堡，將官呂真敗之}；\eventanchor{SYD-1084-XIA-093}{甲辰，夏國主秉常遣使來貢}；\eventanchor{SYD-1084-XIA-098}{六月丙子，夏人寇德順軍，巡檢王友死之}；\eventanchor{SYD-1084-XIA-102}{乙丑，夏人圍定西城，熙河將秦貴敗之}；\eventanchor{SYD-1084-XIA-106}{乙未，夏人寇靜邊砦，涇原將彭孫敗之}。

\subsection{1085年：宋与辽、宋与西夏}
《宋史》、《辽史》、《金史》为这一年留下了可回查的年序入口。

\noindent\textbf{宋与辽。}\eventanchor{SYD-1085-LIAO-070}{癸酉，遣使賀遼主生辰、正旦}；\eventanchor{SYD-1085-LIAO-136}{辛巳，遣使以先帝遺留物遺遼國及告即位，甲申，水部員外郎王諤非職言事，坐罰金}。

\noindent\textbf{宋与西夏。}\eventanchor{SYD-1085-XIA-034}{丙寅，夏人以其母遺留物、馬、白駝來獻}；\eventanchor{SYD-1085-XIA-065}{冬十月甲子，夏國遣使進助山陵馬}；\eventanchor{SYD-1085-XIA-083}{以夏國主母卒，遣使弔祭}。

\subsection{1086年：宋与西夏}
《宋史》、《辽史》、《金史》为这一年留下了可回查的年序入口。

\noindent\textbf{宋与西夏。}\eventanchor{SYD-1086-XIA-035}{庚午，禁邊民與夏人爲市}；\eventanchor{SYD-1086-XIA-066}{庚午，夏國遣使賀坤成節}。

\subsection{1087年：宋、辽与西夏（年序桥接）、宋与西夏}
《宋史》、《辽史》、《续资治通鉴长编》、《金史》为这一年留下了可回查的年序入口。

\noindent\textbf{宋、辽与西夏（年序桥接）。}\eventanchor{SY-1087-CHRONOLOGY-BRIDGE}{【C级年序桥接】保留1087 年的多政权并行检索入口；本条不主张所列政权在当年发生直接互动，具体事件须另以单项材料入账}。

\noindent\textbf{宋与西夏。}\eventanchor{SYD-1087-XIA-036}{二年春正月乙丑，封秉常子乾順爲夏國主}；\eventanchor{SYD-1087-XIA-067}{癸巳，以夏國政亂主幼，強臣乙逋等擅權逆命，詔諸路帥臣嚴兵備之}；\eventanchor{SYD-1087-XIA-084}{辛-\{丑\}-，涇原言夏人寇三川諸砦，官軍敗之}。

\subsection{1088年：宋、辽与西夏（年序桥接）、宋与西夏}
《宋史》、《辽史》、《续资治通鉴长编》、《金史》为这一年留下了可回查的年序入口。

\noindent\textbf{宋、辽与西夏（年序桥接）。}\eventanchor{SY-1088-CHRONOLOGY-BRIDGE}{【C级年序桥接】保留1088 年的多政权并行检索入口；本条不主张所列政权在当年发生直接互动，具体事件须另以单项材料入账}。

\noindent\textbf{宋与西夏。}\eventanchor{SYD-1088-XIA-037}{辛-\{丑\}-，夏人寇塞門砦}；\eventanchor{SYD-1088-XIA-068}{乙亥，夏人寇德靜砦，將官張誠等敗之}。

\subsection{1089年：宋、辽与西夏（年序桥接）、宋与辽、宋与西夏}
《宋史》、《辽史》、《续资治通鉴长编》、《金史》为这一年留下了可回查的年序入口。

\noindent\textbf{宋、辽与西夏（年序桥接）。}\eventanchor{SY-1089-CHRONOLOGY-BRIDGE}{【C级年序桥接】保留1089 年的多政权并行检索入口；本条不主张所列政权在当年发生直接互动，具体事件须另以单项材料入账}。

\noindent\textbf{宋与辽。}\eventanchor{SYD-1089-LIAO-071}{丁-\{丑\}-，遼國使蕭寅等來賀坤成節，曲宴垂拱殿}；\eventanchor{SYD-1089-LIAO-137}{四年春正月壬申朔，不受朝，群臣及遼使詣東上閣門、內東門拜表賀}。

\noindent\textbf{宋与西夏。}\eventanchor{SYD-1089-XIA-038}{甲申，以夏人通好，詔邊將毋生事}；\eventanchor{SYD-1089-XIA-069}{是歲，夏國、邈黎、大食、麻囉拔國入貢}。

\subsection{1090年：宋、辽与西夏（年序桥接）、宋与西夏}
《宋史》、《辽史》、《续资治通鉴长编》、《金史》为这一年留下了可回查的年序入口。

\noindent\textbf{宋、辽与西夏（年序桥接）。}\eventanchor{SY-1090-CHRONOLOGY-BRIDGE}{【C级年序桥接】保留1090 年的多政权并行检索入口；本条不主张所列政权在当年发生直接互动，具体事件须另以单项材料入账}。

\noindent\textbf{宋与西夏。}\eventanchor{SYD-1090-XIA-039}{乙酉，夏人來議分畫疆界}。

\subsection{1091年：宋、辽与西夏（年序桥接）、宋与辽、宋与西夏}
《宋史》、《辽史》、《续资治通鉴长编》、《金史》为这一年留下了可回查的年序入口。

\noindent\textbf{宋、辽与西夏（年序桥接）。}\eventanchor{SY-1091-CHRONOLOGY-BRIDGE}{【C级年序桥接】保留1091 年的多政权并行检索入口；本条不主张所列政权在当年发生直接互动，具体事件须另以单项材料入账}。

\noindent\textbf{宋与辽。}\eventanchor{SYD-1091-LIAO-072}{六年春正月辛酉朔，不受朝，群臣及遼使詣東上閣門、內東門拜表賀}。

\noindent\textbf{宋与西夏。}\eventanchor{SYD-1091-XIA-040}{癸-\{丑\}-，詔鄜延路都監李儀等以違旨夜出兵入界，與夏人戰死，不贈官，餘官降等}；\eventanchor{SYD-1091-XIA-070}{九月丁亥，夏人寇麟、府二州}；\eventanchor{SYD-1091-XIA-085}{夏人寇熙河蘭岷、鄜延路}。

\subsection{1092年：宋与辽}
《宋史》、《辽史》为这一年留下了可回查的年序入口。

\noindent\textbf{宋与辽。}\eventanchor{SYD-1092-LIAO-073}{七年春正月甲辰，以遼使耶律迪卒，輟朝一日}。

\subsection{1093年：宋与辽、宋与西夏}
《宋史》、《辽史》、《金史》为这一年留下了可回查的年序入口。

\noindent\textbf{宋与辽。}\eventanchor{SYD-1093-LIAO-074}{丁巳，遼人遣使來弔祭}。

\noindent\textbf{宋与西夏。}\eventanchor{SYD-1093-XIA-041}{夏四月丁未朔，夏人來謝罪，願以蘭州易塞門砦，不許}。

\subsection{1095年：宋、辽与西夏（年序桥接）}
《宋史》、《辽史》、《续资治通鉴长编》为这一年留下了可回查的年序入口。

\noindent\textbf{宋、辽与西夏（年序桥接）。}\eventanchor{SY-1095-CHRONOLOGY-BRIDGE}{【C级年序桥接】保留1095 年的多政权并行检索入口；本条不主张所列政权在当年发生直接互动，具体事件须另以单项材料入账}。

\subsection{1096年：宋、辽与西夏（年序桥接）、宋与西夏}
《宋史》、《辽史》、《续资治通鉴长编》、《金史》为这一年留下了可回查的年序入口。

\noindent\textbf{宋、辽与西夏（年序桥接）。}\eventanchor{SY-1096-CHRONOLOGY-BRIDGE}{【C级年序桥接】保留1096 年的多政权并行检索入口；本条不主张所列政权在当年发生直接互动，具体事件须另以单项材料入账}。

\noindent\textbf{宋与西夏。}\eventanchor{SYD-1096-XIA-042}{八月辛酉，夏人寇寧順砦}；\eventanchor{SYD-1096-XIA-071}{壬戌，夏人寇鄜、延，陷金明砦}。

\subsection{1097年：宋与西夏}
《宋史》、《辽史》、《金史》为这一年留下了可回查的年序入口。

\noindent\textbf{宋与西夏。}\eventanchor{SYD-1097-XIA-043}{庚午，夏人大至葭盧城下，知石州張構等擊走之}；\eventanchor{SYD-1097-XIA-072}{庚子，知保安軍李沂伐夏國，破洪州}；\eventanchor{SYD-1097-XIA-086}{癸亥，于闐來貢，黑汗王攻夏人三州，遣其子以聞}；\eventanchor{SYD-1097-XIA-094}{甲午，涇原路鈐轄王文振敗夏人於沒煙峽}；\eventanchor{SYD-1097-XIA-099}{壬寅，環慶鈐轄張存入鹽州，俘戮甚衆，及還，夏人追襲之，復多亡失}；\eventanchor{SYD-1097-XIA-103}{三月壬戌，夏人犯麟州神堂堡，出兵討之，及進築胡山砦}；\eventanchor{SYD-1097-XIA-107}{辛巳，西上閣門使折克行破夏人于長波川，斬首二千餘級，獲牛馬倍之}。

\subsection{1098年：宋、辽与西夏（年序桥接）、宋与西夏}
《宋史》、《辽史》、《续资治通鉴长编》、《金史》为这一年留下了可回查的年序入口。

\noindent\textbf{宋、辽与西夏（年序桥接）。}\eventanchor{SY-1098-CHRONOLOGY-BRIDGE}{【C级年序桥接】保留1098 年的多政权并行检索入口；本条不主张所列政权在当年发生直接互动，具体事件须另以单项材料入账}。

\noindent\textbf{宋与西夏。}\eventanchor{SYD-1098-XIA-044}{涇原路禽夏國統軍嵬名阿埋等，高麗、瞎征、西南蕃張氏、羅氏、程氏入貢}。

\subsection{1099年：宋、辽与西夏（年序桥接）、宋与辽、宋与西夏}
《宋史》、《辽史》、《续资治通鉴长编》、《金史》为这一年留下了可回查的年序入口。

\noindent\textbf{宋、辽与西夏（年序桥接）。}\eventanchor{SY-1099-CHRONOLOGY-BRIDGE}{【C级年序桥接】保留1099 年的多政权并行检索入口；本条不主张所列政权在当年发生直接互动，具体事件须另以单项材料入账}。

\noindent\textbf{宋与辽。}\eventanchor{SYD-1099-LIAO-075}{二年春正月甲辰朔，-\{御\}-大慶殿，以雪罷朝，群臣及遼使詣東上閣門拜表賀}；\eventanchor{SYD-1099-LIAO-138}{遣中書舍人郭知章報聘於遼}。

\noindent\textbf{宋与西夏。}\eventanchor{SYD-1099-XIA-045}{甲申，夏人以國母卒，遣使告哀，且謝罪，卻其使不納}；\eventanchor{SYD-1099-XIA-073}{九月庚子朔，夏人來謝罪}；\eventanchor{SYD-1099-XIA-087}{戊子，鄜延鈐轄劉安敗夏人於神堆}。

\subsection{1100年：宋与辽}
《宋史》、《辽史》为这一年留下了可回查的年序入口。

\noindent\textbf{宋与辽。}\eventanchor{SYD-1100-LIAO-076}{丙午，遣董敦逸賀遼主生辰，呂仲甫賀正旦}；\eventanchor{SYD-1100-LIAO-139}{丙寅，遼遣蕭穆來賀即位}。

\subsection{1101年：宋与辽}
《宋史》、《辽史》为这一年留下了可回查的年序入口。

\noindent\textbf{宋与辽。}\eventanchor{SYD-1101-LIAO-077}{是歲，遼人來獻遺留物}；\eventanchor{SYD-1101-LIAO-140}{乙丑，遼使蕭恭來告其主洪基殂，遣謝瓘、上官均等往弔祭，黃寔賀其孫延禧立}。

\subsection{1104年：宋与西夏}
《宋史》、《辽史》、《金史》为这一年留下了可回查的年序入口。

\noindent\textbf{宋与西夏。}\eventanchor{SYD-1104-XIA-046}{戊午，夏人入涇原，圍平夏城，寇鎮戎軍}。

\subsection{1105年：宋、辽与西夏（年序桥接）、宋与辽、宋与西夏}
《宋史》、《辽史》、《续资治通鉴长编》、《金史》为这一年留下了可回查的年序入口。

\noindent\textbf{宋、辽与西夏（年序桥接）。}\eventanchor{SY-1105-CHRONOLOGY-BRIDGE}{【C级年序桥接】保留1105 年的多政权并行检索入口；本条不主张所列政权在当年发生直接互动，具体事件须另以单项材料入账}。

\noindent\textbf{宋与辽。}\eventanchor{SYD-1105-LIAO-078}{壬子，遣林攄報聘於遼}。

\noindent\textbf{宋与西夏。}\eventanchor{SYD-1105-XIA-047}{己丑，夏人寇順寧砦，鄜延第二副將劉延慶擊破之}。

\subsection{1108年：宋、辽与西夏（年序桥接）}
《宋史》、《辽史》、《续资治通鉴长编》为这一年留下了可回查的年序入口。

\noindent\textbf{宋、辽与西夏（年序桥接）。}\eventanchor{SY-1108-CHRONOLOGY-BRIDGE}{【C级年序桥接】保留1108 年的多政权并行检索入口；本条不主张所列政权在当年发生直接互动，具体事件须另以单项材料入账}。

\subsection{1109年：宋、辽与西夏（年序桥接）、宋与西夏}
《宋史》、《辽史》、《续资治通鉴长编》、《金史》为这一年留下了可回查的年序入口。

\noindent\textbf{宋、辽与西夏（年序桥接）。}\eventanchor{SY-1109-CHRONOLOGY-BRIDGE}{【C级年序桥接】保留1109 年的多政权并行检索入口；本条不主张所列政权在当年发生直接互动，具体事件须另以单项材料入账}。

\noindent\textbf{宋与西夏。}\eventanchor{SYD-1109-XIA-048}{闍婆、占城、夏國入貢}。

\subsection{1110年：宋、辽与西夏（年序桥接）}
《宋史》、《辽史》、《续资治通鉴长编》为这一年留下了可回查的年序入口。

\noindent\textbf{宋、辽与西夏（年序桥接）。}\eventanchor{SY-1110-CHRONOLOGY-BRIDGE}{【C级年序桥接】保留1110 年的多政权并行检索入口；本条不主张所列政权在当年发生直接互动，具体事件须另以单项材料入账}。

\subsection{1111年：宋、辽与西夏（年序桥接）、宋与辽}
《宋史》、《辽史》、《续资治通鉴长编》为这一年留下了可回查的年序入口。

\noindent\textbf{宋、辽与西夏（年序桥接）。}\eventanchor{SY-1111-CHRONOLOGY-BRIDGE}{【C级年序桥接】保留1111 年的多政权并行检索入口；本条不主张所列政权在当年发生直接互动，具体事件须另以单项材料入账}。

\noindent\textbf{宋与辽。}\eventanchor{SYD-1111-LIAO-079}{是月，鄭允中、童貫使遼，以李良嗣來，良嗣獻取燕之策，詔賜姓趙}。

\subsection{1112年：宋、辽与西夏（年序桥接）}
《宋史》、《辽史》、《续资治通鉴长编》为这一年留下了可回查的年序入口。

\noindent\textbf{宋、辽与西夏（年序桥接）。}\eventanchor{SY-1112-CHRONOLOGY-BRIDGE}{【C级年序桥接】保留1112 年的多政权并行检索入口；本条不主张所列政权在当年发生直接互动，具体事件须另以单项材料入账}。

\subsection{1113年：宋、辽与西夏（年序桥接）、宋与辽}
《宋史》、《辽史》、《续资治通鉴长编》为这一年留下了可回查的年序入口。

\noindent\textbf{宋、辽与西夏（年序桥接）。}\eventanchor{SY-1113-CHRONOLOGY-BRIDGE}{【C级年序桥接】保留1113 年的多政权并行检索入口；本条不主张所列政权在当年发生直接互动，具体事件须另以单项材料入账}。

\noindent\textbf{宋与辽。}\eventanchor{SYD-1113-LIAO-080}{甲午，以遼、女真相持，詔河北治邊防}。

\subsection{1116年：宋、辽与金（年序桥接）、宋与西夏}
《宋史》、《辽史》、《金史》、《续资治通鉴长编》为这一年留下了可回查的年序入口。

\noindent\textbf{宋、辽与金（年序桥接）。}\eventanchor{SY-1116-CHRONOLOGY-BRIDGE}{【C级年序桥接】保留1116 年的多政权并行检索入口；本条不主张所列政权在当年发生直接互动，具体事件须另以单项材料入账}。

\noindent\textbf{宋与西夏。}\eventanchor{SYD-1116-XIA-049}{高麗、占城、大食、真臘、大理、夏國入貢，茂州夷郅永壽內附}。

\subsection{1118年：宋与辽}
《宋史》、《辽史》为这一年留下了可回查的年序入口。

\noindent\textbf{宋与辽。}\eventanchor{SYD-1118-LIAO-081}{庚午，遣武義大夫馬政由海道使女真，約夾攻遼}。

\subsection{1119年：宋与西夏}
《宋史》、《辽史》、《金史》为这一年留下了可回查的年序入口。

\noindent\textbf{宋与西夏。}\eventanchor{SYD-1119-XIA-050}{庚寅，童貫以鄜延、環慶兵大破夏人，平其三城}；\eventanchor{SYD-1119-XIA-074}{己亥，夏國遣使納款，詔六路罷兵}；\eventanchor{SYD-1119-XIA-088}{童貫遣知熙州劉法出師攻統安城，夏人伏兵擊之，法敗歿，震武軍受圍}。

\subsection{1122年：宋与辽}
《宋史》、《辽史》为这一年留下了可回查的年序入口。

\noindent\textbf{宋与辽。}\eventanchor{SYD-1122-LIAO-082}{丙子，遼人立燕王淳為帝}。

\subsection{1123年：宋与辽}
《宋史》、《辽史》为这一年留下了可回查的年序入口。

\noindent\textbf{宋与辽。}\eventanchor{SYD-1123-LIAO-083}{丙戌，遼人張覺以平州來附}。

\subsection{1124年：宋与西夏}
《宋史》、《辽史》、《金史》为这一年留下了可回查的年序入口。

\noindent\textbf{宋与西夏。}\eventanchor{SYD-1124-XIA-051}{夏國、高麗、于闐、羅殿入貢}。

\subsection{1125年：宋与辽}
《宋史》、《辽史》为这一年留下了可回查的年序入口。

\noindent\textbf{宋与辽。}\eventanchor{SYD-1125-LIAO-084}{壬辰，金人以擒遼主，遣李孝和等來告慶}。

\subsection{1126年：宋与辽、宋与西夏}
《宋史》、《辽史》、《金史》为这一年留下了可回查的年序入口。

\noindent\textbf{宋与辽。}\eventanchor{SYD-1126-LIAO-085}{遼故將小䩴䩮攻陷麟州建寧砦，知砦楊震死之}。

\noindent\textbf{宋与西夏。}\eventanchor{SYD-1126-XIA-052}{十一月丙寅，夏人陷懷德軍，知軍事劉銓、通判杜翊世死之}；\eventanchor{SYD-1126-XIA-075}{夏四月戊戌，夏人陷震威城，攝知城事朱昭死之}。
